% Section 7.1 — The Limits of Fixed-Threshold Cataloging
% Combines 7.1_introduction, 7.1.1_catalog_paradigm, 7.1.2_info_below_threshold

\section{The Limits of Fixed-Threshold Cataloging}
\label{sec:threshold_limits}

Before describing how the source-count distribution can be leveraged to build a probabilistic catalog, it is useful to examine more closely the limitations of the standard approach.
The fixed-threshold paradigm that underpins the \textit{Fermi}-LAT catalogs has been used effectively for over a decade, but it involves two distinct costs: a spatially non-uniform detection significance and the wholesale discarding of sub-threshold information.
We address these in turn.

\subsection{The Standard Catalog-Construction Paradigm}
\label{sec:catalog_paradigm}

The \textit{Fermi}-LAT fourth source catalog (4FGL) is the most comprehensive catalog produced from the instrument described in Chapter~\ref{chap:gamma_sky}, incorporating over a decade of survey-mode observations in the gamma-ray band~\cite{Fermi-LAT:2019yla}.
Point-source detection in this catalog follows a well-defined statistical procedure.
In each region of interest (RoI) partitioning the sky, the analysis evaluates a likelihood-ratio test statistic defined as $\mathrm{TS} = 2\ln(\mathcal{L}/\mathcal{L}_0)$, where $\mathcal{L}$ is the maximum likelihood including the source model and $\mathcal{L}_0$ is the maximum likelihood for the background-only hypothesis~\cite{Fermi-LAT:2019yla}. Sources are retained in the catalog if they exceed a threshold of $\mathrm{TS} > 25$, corresponding to a significance of approximately $4\sigma$ when evaluated against a $\chi^2$ distribution with four degrees of freedom (two positional and two spectral parameters for a power-law model)~\cite{Fermi-LAT:2019yla}.
Any candidate that falls below this threshold is removed from the model, and the fit in that RoI is readjusted accordingly.

This detection procedure is inherently local.
The sky is divided into roughly 1750 RoIs, each containing a core region of actively fitted sources and an outer ring whose width adapts to the energy-dependent point-spread function~\cite{Fermi-LAT:2019yla}. Within each RoI, the normalization and spectral tilt of the Galactic diffuse emission template and the normalization of the isotropic component are fitted independently, and the log-likelihood is further split into components defined by energy intervals and event types before being summed for optimization~\cite{Fermi-LAT:2019yla}.
Since the background model parameters are determined locally, the same numerical TS value carries different statistical weight at different sky positions.
A source yielding $\mathrm{TS} = 30$ in a sparse high-latitude field sits on a very different background than one with $\mathrm{TS} = 30$ in the crowded Galactic plane.
The resulting detection scale is therefore spatially inhomogeneous: the TS encodes the local excess above the locally fitted background, not a globally calibrated measure of source significance.

The fixed detection threshold imposes a binary classification on the sky.
Every direction is either ``in'' the catalog ($\mathrm{TS} \geq 25$) or ``out,'' and the information carried by pixels with $\mathrm{TS}$ values just below the cut --- pixels that may well host genuine astrophysical sources --- is entirely discarded.
The construction of each successive data release inherits this binary structure.
While such a threshold is operationally necessary to maintain a controlled false-positive rate across thousands of candidate seeds, it comes at the cost of throwing away a continuous gradient of statistical evidence that, as we will see, can be recovered through comparison with simulated sky ensembles.

\subsection{The Information Below Threshold}
\label{sec:info_below_threshold}

The fixed detection threshold does not merely exclude noise.
A large population of genuine astrophysical sources lies just below $\mathrm{TS} = 25$, and the collective emission of these objects constitutes a significant fraction of the observed high-latitude gamma-ray sky.
Blazars, misaligned active galactic nuclei, star-forming galaxies, and millisecond pulsars are all guaranteed contributors to this unresolved emission~\cite{Fornasa:2015qua}.
Of these, blazars dominate the resolved source counts but account for roughly twenty percent of the unresolved gamma-ray background in the $0.1$--$100~\mathrm{GeV}$ range; misaligned active galactic nuclei contribute a comparable fraction ($\sim\!
	25\%$), while star-forming galaxies may account for between four and fifty percent depending on the assumed luminosity function and spectral model~\cite{Fornasa:2015qua}. Millisecond pulsars are subdominant, responsible for less than one percent of the unresolved intensity~\cite{Fornasa:2015qua}. Taken together, these populations can explain the measured intensity spectrum of the isotropic diffuse background without exotic contributions~\cite{Fornasa:2015qua, Fermi-LAT:2014ryh}, though the individual budget shares remain uncertain.

The source-count distribution recovered in Chapter~\ref{ch:6} (cf.\ Section~\ref{sec:source_count}) quantifies this sub-threshold population directly. The $dN/dS$ measured from the high-latitude \textit{Fermi}-LAT photon counts extends as $\sim S^{-2}$ down to fluxes of approximately $5 \times 10^{-12}~\mathrm{cm}^{-2}\,\mathrm{s}^{-1}$, roughly 50 times below the nominal catalog sensitivity of $\sim 2 \times 10^{-10}~\mathrm{cm}^{-2}\,\mathrm{s}^{-1}$ in the $1$--$10~\mathrm{GeV}$ band~\cite{Korsmeier:2022cwp} (cf.\ the analysis presented in Chapter~\ref{ch:6}).
This faint extension implies a substantial reservoir of sources whose individual signals fall just below the $\mathrm{TS} = 25$ cut-off --- sources that carry real astrophysical information yet are systematically excluded from every standard catalog.
For science cases that depend on sample size rather than per-source purity, this information loss is particularly costly.
Cross-correlation studies between gamma-ray maps and galaxy catalogs, for instance, gain statistical power from a larger sample of source directions, even at the price of a moderate false-positive fraction (cf.
\ Chapter~\ref{ch:8}).
Similarly, multi-wavelength source identification campaigns benefit from having a larger set of candidate positions to follow up.

The practical consequences of the fixed threshold become visible in the incremental catalog releases.
In the 4FGL-DR3, which extends the observation baseline from eight to twelve years, 112 sources originally included in the first data release dropped below $\mathrm{TS} = 25$, fifteen of which fell below $\mathrm{TS} = 16$, with the lowest reaching $\mathrm{TS} = 9.8$~\cite{Fermi-LAT:2022byn}. An additional 69 sources from the second data release also fell below threshold~\cite{Fermi-LAT:2022byn}. These sources were retained in the catalog for continuity and explicitly flagged, but the episode illustrates a structural difficulty: highly variable objects such as flat-spectrum radio quasars can be bright enough to enter the catalog during one observation window and then fade below the detection threshold in subsequent data accumulations~\cite{Bhat:2021wtb}. Simply lowering the threshold does not resolve this problem, since it would increase the number of false detections from background fluctuations~\cite{Bhat:2021wtb}.What is needed instead is a detection framework that replaces the binary in-or-out decision with a continuous, simulation-calibrated measure of source probability --- one that accounts for the expected sub-threshold population as derived from the $dN/dS$.
